\chapter{设计模式与工程实践总结}

\section{架构级模式}

\subsection{并行预取 (Parallel Prefetch)}

Claude Code 最具特色的架构模式是\textbf{利用模块加载时间进行并行预取}：

\begin{lstlisting}[language=TypeScript,breaklines=true]
// 在 import 语句之前启动异步 I/O
profileCheckpoint('main_tsx_entry');
startMdmRawRead();        // 启动 MDM 子进程
startKeychainPrefetch();  // 启动钥匙串读取

// 后续的 import 语句需要 ~135ms
// 此时 MDM 和 Keychain 读取在后台进行
import chalk from 'chalk';
import React from 'react';
// ...
\end{lstlisting}

\textbf{原则}：识别"免费"的等待时间，将 I/O 操作前移到该窗口。

\subsection{分层延迟初始化}

\begin{lstlisting}[language=TypeScript,breaklines=true]
阶段 1：模块求值前 → 启动异步 I/O
阶段 2：模块求值中 → 自动完成（JS 引擎）
阶段 3：首次渲染前 → 必要的同步初始化
阶段 4：首次渲染后 → 延迟预取（用户不感知）
\end{lstlisting}

\subsection{AsyncGenerator 驱动的流式处理}

\begin{lstlisting}[language=TypeScript,breaklines=true]
async function* query(params): AsyncGenerator<StreamEvent | Message> {
  while (true) {
    yield { type: 'stream_request_start' }
    // ... API 调用 & 工具执行 ...
    yield assistantMessage
    // ... 判断是否继续 ...
  }
}
\end{lstlisting}

使用 AsyncGenerator 实现流式数据管道，允许调用方实时消费每个事件。

\section{代码组织模式}

\subsection{条件导入与死代码消除}

\begin{lstlisting}[language=TypeScript,breaklines=true]
const SleepTool = feature('PROACTIVE')
  ? require('./tools/SleepTool/SleepTool.js').SleepTool
  : null
\end{lstlisting}

\textbf{规律}：所有可选功能都使用三元表达式 + \code{require()} 而非顶层 \code{import}。这使得 Bun 的 bundler 可以在编译时完全移除未使用的模块。

\subsection{懒加载打破循环依赖}

\begin{lstlisting}[language=TypeScript,breaklines=true]
// 不使用顶层 import（会创建循环依赖）
const getTeamCreateTool = () =>
  require('./tools/TeamCreateTool/TeamCreateTool.js').TeamCreateTool
\end{lstlisting}

\textbf{规律}：将导入包装在函数中，只在实际调用时执行。搭配 TypeScript 的 \code{typeof import(...)} 保留类型安全。

\subsection{Memoize 模式}

\begin{lstlisting}[language=TypeScript,breaklines=true]
export const getSystemContext = memoize(async () => {
  // 只执行一次的重量级操作
  const gitStatus = await getGitStatus()
  return { gitStatus }
})
\end{lstlisting}

对于整个会话期间不变的数据（Git 状态、CLAUDE.md 内容），使用 \code{memoize} 缓存结果。

\subsection{清除缓存的信号}

\begin{lstlisting}[language=TypeScript,breaklines=true]
export function setSystemPromptInjection(value: string | null): void {
  systemPromptInjection = value
  getUserContext.cache.clear?.()    // 清除 memoize 缓存
  getSystemContext.cache.clear?.()
}
\end{lstlisting}

当底层数据变化时，手动清除 memoize 缓存。

\section{类型系统模式}

\subsection{DeepImmutable}

\begin{lstlisting}[language=TypeScript,breaklines=true]
export type ToolPermissionContext = DeepImmutable<{
  mode: PermissionMode
  alwaysAllowRules: ToolPermissionRulesBySource
  // ...
}>
\end{lstlisting}

使用 \code{DeepImmutable} 确保权限上下文在传递过程中不被意外修改。

\subsection{自文档化类型名}

\begin{lstlisting}[language=TypeScript,breaklines=true]
type AnalyticsMetadata_I_VERIFIED_THIS_IS_NOT_CODE_OR_FILEPATHS = string
\end{lstlisting}

类型名本身就是一条检查清单。每次使用时，开发者必须 pause 并确认。

\subsection{品牌类型 (Branded Types)}

\begin{lstlisting}[language=TypeScript,breaklines=true]
import { asSessionId } from './types/ids.js'
import type { AgentId } from './types/ids.js'
\end{lstlisting}

使用品牌类型防止 ID 类型之间的混用。

\section{错误处理模式}

\subsection{恢复循环}

\begin{lstlisting}[language=TypeScript,breaklines=true]
const MAX_OUTPUT_TOKENS_RECOVERY_LIMIT = 3

// 自动恢复，最多重试 3 次
// 恢复期间扣留错误消息，避免影响 SDK 调用方
\end{lstlisting}

\subsection{优雅降级}

\begin{lstlisting}[language=TypeScript,breaklines=true]
const truncatedStatus = status.length > MAX_STATUS_CHARS
  ? status.substring(0, MAX_STATUS_CHARS) +
    '\n... (truncated. Use BashTool for full info)'
  : status
\end{lstlisting}

超出限制时截断并提供获取完整信息的替代方案。

\subsection{分类重试}

\begin{lstlisting}[language=TypeScript,breaklines=true]
import { categorizeRetryableAPIError } from './services/api/errors.js'
import { FallbackTriggeredError } from './services/api/withRetry.js'
\end{lstlisting}

API 错误被分类为可重试和不可重试，支持自动重试和模型 fallback。

\section{安全模式}

\subsection{纵深防御}

\begin{lstlisting}[language=TypeScript,breaklines=true]
1. 工具列表过滤 → 不暴露被拒工具
2. 调用时权限检查 → 每次操作验证
3. 沙箱执行 → 受限环境
4. 审计追踪 → 记录所有决策
\end{lstlisting}

\subsection{隐私优先}

\begin{lstlisting}[language=TypeScript,breaklines=true]
logForDiagnosticsNoPII(...)     // 日志不含 PII
_I_VERIFIED_THIS_IS_NOT_CODE_OR_FILEPATHS  // 类型级别隐私检查
\end{lstlisting}

\subsection{信任边界}

\begin{lstlisting}[language=TypeScript,breaklines=true]
if (!hasTrust) {
  // 不执行 Git 命令（可能通过 hooks 执行任意代码）
}
\end{lstlisting}

\section{性能模式}

\subsection{并行化}

\begin{lstlisting}[language=TypeScript,breaklines=true]
// 5 个 Git 命令并行
const [branch, mainBranch, status, log, userName] = await Promise.all([...])

// 启动时并行预取
startMdmRawRead()
startKeychainPrefetch()
void getSystemContext()
\end{lstlisting}

\subsection{缓存感知排序}

\begin{lstlisting}[language=TypeScript,breaklines=true]
// 内置工具排在前面，保持稳定排序
// 使 Prompt Cache 跨请求复用
return uniqBy(
  [...builtInTools].sort(byName).concat(allowedMcpTools.sort(byName)),
  'name',
)
\end{lstlisting}

\subsection{延迟加载重量级模块}

\begin{lstlisting}[language=TypeScript,breaklines=true]
OpenTelemetry (~400KB) → 动态 import()
gRPC (~700KB) → 动态 import()
\end{lstlisting}

\section{项目规模管理}

一个 512,000+ 行的 TypeScript 项目面临的挑战和应对：

\begin{longtable}{ll}
\toprule
挑战 & Claude Code 的应对 \\
\midrule
循环依赖 & 懒加载 \code{require()} 函数 \\
启动速度 & 并行预取 + 延迟初始化 \\
Bundle 大小 & Feature Flag + 死代码消除 \\
类型安全 & 严格模式 + 品牌类型 \\
配置演进 & 版本化迁移系统 \\
隐私合规 & 类型级别隐私标记 \\
\bottomrule
\end{longtable}

\section{总结}

Claude Code 的源码揭示了构建工业级 AI 编程助手的复杂度。它不仅仅是一个"调用 API 的 CLI 工具"，而是一个包含以下系统的完整平台：

\begin{enumerate}[nosep]
  \item \textbf{查询引擎}：多层压缩管线、自动恢复、Token 预算
  \item \textbf{工具系统}：40+ 工具、权限模型、缓存感知
  \item \textbf{命令系统}：50+ 命令、生命周期管理
  \item \textbf{协调器}：多代理编排、并行执行
  \item \textbf{UI 框架}：React + Ink 的 140+ 组件
  \item \textbf{安全层}：多模式权限、沙箱、审计
  \item \textbf{扩展系统}：插件、技能、MCP
  \item \textbf{遥测}：OpenTelemetry、成本追踪、性能分析
\end{enumerate}

这些系统的设计模式——并行预取、缓存感知、分层压缩、纵深防御——不仅适用于 AI 编程助手，也可以推广到任何需要将 LLM 深度集成到软件工程流程中的应用。

\bigskip\noindent\rule{\textwidth}{0.4pt}\bigskip

\textbf{全书完}

> 本书基于 \code{szylover/claude-code} 仓库中的公开源码撰写。
> 所有原始源码的知识产权归 \href{https://www.anthropic.com}{Anthropic} 所有。
